% Provenance (never printed): NSI terminale "Architectures materielles, systemes d'exploitation et reseaux: gestion
% des processus et des ressources par un systeme d'exploitation (creation, ordonnancement, interblocage); observer
% les processus actifs ou en attente avec une commande" (+ interrupts, named in the introduction). See
% sources/book2/BRIEFS.md.
\chapter{Processes and Scheduling}\label{ch:g12:processes-and-scheduling}

A laptop with four processor cores runs some three hundred programs ``at the same time'': the editor, the browser and
its tabs, the music player, and many small programs of the operating system that nobody started by hand. Four cores
can run four instructions at once, not three hundred. The rest is an illusion that the operating system maintains by
switching from one program to the next many times a second.

This chapter looks at a running program, a \omterm{def:g12:processes-and-scheduling:process}{process}, from the operating system's side: how it is created, the \omterm{def:g12:processes-and-scheduling:states}{states}
it passes through, and how the system chooses which one runs. Sharing the processor raises two new kinds of bug, a
lost update and a \omterm{def:g12:processes-and-scheduling:deadlock}{deadlock}, which appear only when programs run together; both are shown on small simulations where
every possible order can be tried.

\begin{recall}
The \omterm{def:g11:os-and-command-line:kernel}{kernel} is the part of the operating system that manages the hardware and runs programs; programs ask it for
services by \omterm{def:g11:os-and-command-line:kernel}{system calls}, and the shell starts commands (chapter \ref{ch:g11:os-and-command-line}). A \omterm{def:g10:sensors-actuators-control:state-machine}{state machine}
moves between \omterm{def:g12:processes-and-scheduling:states}{states} on events (chapter \ref{ch:g10:sensors-actuators-control}). A queue serves its elements first in,
first out (chapter \ref{ch:g12:lists-stacks-queues}). A cycle in a \omterm{def:g10:graphs-first-look:directed}{directed graph} is a path that comes back to its
start (chapter \ref{ch:g12:graph-algorithms}).
\end{recall}

\section{Programs and processes}\label{sec:g12:processes-and-scheduling:process}

\begin{definition}[Process]\label{def:g12:processes-and-scheduling:process}
A \emph{process}\index{process} is a program being run. It has its own memory (its code, its data, its stack), the
values of the processor's registers when it last ran, the files it has opened, and a number, its process identifier,
given by the operating system.
\end{definition}

One program can give several processes: two windows of the same editor are two processes running the same code with
different data. A \omterm{def:g12:processes-and-scheduling:process}{process} creates another by a \omterm{def:g11:os-and-command-line:kernel}{system call}; the new one is its \emph{child}. In Python, the module
\code{subprocess} starts a child and waits for it.

\omcode{code/grade-12/10-processes-and-scheduling/listings/spawn.py}{1}{10}{A parent \omterm{def:g12:processes-and-scheduling:process}{process} starts two
children.}\label{lst:g12:processes-and-scheduling:spawn}

\omoutput{code/grade-12/10-processes-and-scheduling/outputs/spawn.txt}{Each child ends with an exit status: $0$ means
success, anything else a failure that the parent can test.}\label{out:g12:processes-and-scheduling:spawn}

The command \code{ps} lists processes. Below, a shell starts \code{sleep 30} in the background with \texttt{\&}, so that
the shell does not wait for it; \texttt{\$!} is the number of that last background \omterm{def:g12:processes-and-scheduling:process}{process}, and \texttt{\$\$} the
shell's own. The option \code{-p} selects processes by number, and \code{-o state=,comm=} asks for two columns: the
state and the command.
% ledger: g12:ps-state

\omoutput{code/grade-12/10-processes-and-scheduling/outputs/ps_session.txt}{Both processes are in state \code{S},
asleep: \code{sleep} waits for its thirty seconds to pass, the shell for \code{ps} to
finish.}\label{out:g12:processes-and-scheduling:ps}

\section{The states of a process}\label{sec:g12:processes-and-scheduling:states}

\begin{definition}[Process states]\label{def:g12:processes-and-scheduling:states}
At each moment a \omterm{def:g12:processes-and-scheduling:process}{process} is in one of three \emph{states}\index{process state (ready, running, blocked)}:
\emph{running}, if a processor core is executing it; \emph{ready}, if it could run but waits for a core;
\emph{blocked}, if it waits for an event, such as the end of a disk read, a key press or a message from the network,
and could not use a core if it had one.
\end{definition}

\begin{omfigure}
\begin{tikzpicture}[font=\footnotesize, st/.style={draw, circle, minimum size=16mm, align=center, fill=omDef!10},
    ed/.style={->, thick, >={Stealth}}]
  \node[st] (rd) at (0,0) {ready};
  \node[st, fill=omMeth!15] (rn) at (5,0) {running};
  \node[st, fill=omThm!12] (bl) at (2.5,-3.6) {blocked};
  \node[draw, rounded corners, inner sep=3pt] (nw) at (-3,0) {created};
  \node[draw, rounded corners, inner sep=3pt] (dn) at (8,0) {ended};
  \draw[ed] (nw) -- (rd);
  \draw[ed] (rd) to[bend left=25] node[above] {chosen by the scheduler} (rn);
  \draw[ed] (rn) to[bend left=25] node[below, align=center] {interrupted:\\ quantum over} (rd);
  \draw[ed] (rn) -- node[pos=0.6, right=2pt, align=left] {waits for\\ an event} (bl);
  \draw[ed] (bl) -- node[pos=0.4, left=2pt, align=right] {the event\\ happened} (rd);
  \draw[ed] (rn) -- node[above] {exits} (dn);
\end{tikzpicture}

{\small The \omterm{def:g12:processes-and-scheduling:states}{states} of a \omterm{def:g12:processes-and-scheduling:process}{process} as a \omterm{def:g10:sensors-actuators-control:state-machine}{state machine}. Only the \omterm{def:g12:processes-and-scheduling:scheduler}{scheduler} moves a \omterm{def:g12:processes-and-scheduling:process}{process} into \emph{running}; a blocked
\omterm{def:g12:processes-and-scheduling:process}{process} always goes back through \emph{ready}, even when its event is over.}
\end{omfigure}

The output of \code{ps} above uses letters for these \omterm{def:g12:processes-and-scheduling:states}{states}: \code{R} for running or ready, \code{S} for asleep, waiting for an
event to complete.
% ledger: g12:ps-state

\section{The scheduler}\label{sec:g12:processes-and-scheduling:scheduler}

\begin{definition}[Interrupt]\label{def:g12:processes-and-scheduling:interrupt}
An \emph{interrupt}\index{interrupt} is a signal from a device to the processor that makes it stop the current
program, at the end of the current instruction, and run a function of the \omterm{def:g11:os-and-command-line:kernel}{kernel}. A timer, the network card or the
disk controller each sends interrupts; the \omterm{def:g11:os-and-command-line:kernel}{kernel}'s function decides what happens next.
\end{definition}

A timer \omterm{def:g12:processes-and-scheduling:interrupt}{interrupts} the processor at a regular pace, many times a second. Each time, the \omterm{def:g11:os-and-command-line:kernel}{kernel} can take the processor away from the running \omterm{def:g12:processes-and-scheduling:process}{process} and give it to another.

\begin{definition}[Scheduler, quantum, round robin]\label{def:g12:processes-and-scheduling:scheduler}
The \emph{scheduler}\index{scheduler} is the part of the \omterm{def:g11:os-and-command-line:kernel}{kernel} that chooses which ready \omterm{def:g12:processes-and-scheduling:process}{process} runs next. A
\emph{quantum}\index{quantum} is the longest time a \omterm{def:g12:processes-and-scheduling:process}{process} may run before it is interrupted and sent back to the
ready queue. \emph{Round robin}\index{round robin} keeps the ready processes in a queue: it runs the first one for at
most a quantum, then puts it back at the end of the queue if it is not finished. \emph{First come, first served} runs
each \omterm{def:g12:processes-and-scheduling:process}{process} to its end, in order of arrival.
\end{definition}

\begin{definition}[Context switch]\label{def:g12:processes-and-scheduling:context-switch}
A \emph{context switch}\index{context switch} is the work of the \omterm{def:g11:os-and-command-line:kernel}{kernel} to stop one \omterm{def:g12:processes-and-scheduling:process}{process} and start another: save
the registers of the first, choose the second, restore its registers and its memory map. While it happens, no \omterm{def:g12:processes-and-scheduling:process}{process}
does useful work.
\end{definition}

The simulator below follows these rules on a \emph{workload}: a list of processes, each with its arrival time and
the processor time it needs, its \emph{burst}. Time is counted in units, the same for all. A \omterm{def:g12:processes-and-scheduling:process}{process}'s
\emph{turnaround} is its finish time minus its arrival; its \emph{waiting time} is its turnaround minus its burst; its
\emph{response time} is the time from its arrival to its first run.

\omcode{code/grade-12/10-processes-and-scheduling/listings/scheduler.py}{1}{39}{First come, first served and \omterm{def:g12:processes-and-scheduling:scheduler}{round
robin}, simulated.}\label{lst:g12:processes-and-scheduling:scheduler}

\omcode{code/grade-12/10-processes-and-scheduling/listings/scheduler.py}{42}{56}{The times of each \omterm{def:g12:processes-and-scheduling:process}{process}, and a
Gantt line of text.}\label{lst:g12:processes-and-scheduling:times}

\omcode{code/grade-12/10-processes-and-scheduling/listings/fcfs_rr.py}{1}{16}{One workload, three
schedules.}\label{lst:g12:processes-and-scheduling:fcfs-rr}

\omoutput{code/grade-12/10-processes-and-scheduling/outputs/fcfs_rr.txt}{Each letter is one unit of processor time
given to that \omterm{def:g12:processes-and-scheduling:process}{process}. \omterm{def:g12:processes-and-scheduling:scheduler}{Round robin} starts every \omterm{def:g12:processes-and-scheduling:process}{process} sooner, at the price of a longer
turnaround.}\label{out:g12:processes-and-scheduling:fcfs-rr}

\begin{omfigure}
\begin{tikzpicture}[font=\footnotesize, x=7mm, y=7mm]
  \foreach \n/\a/\b/\c in {A/0/2/omDef!12, B/2/4/omMeth!18, C/4/6/omThm!14, A/6/8/omDef!12, D/8/10/black!10,
      B/10/11/omMeth!18, C/11/13/omThm!14, A/13/14/omDef!12, C/14/18/omThm!14}
    \draw[fill=\c] (\a,0) rectangle node {\n} (\b,1);
  \foreach \t in {0,2,...,18} \node[below, font=\scriptsize] at (\t,0) {\t};
  \foreach \n/\t in {A/0, B/1, C/2, D/3} \draw[->, omThm] (\t,2) node[above, font=\scriptsize] {\n} -- (\t,1.05);
  \node[left, font=\scriptsize] at (-0.3,1.6) {arrivals};
\end{tikzpicture}

{\small \omterm{def:g12:processes-and-scheduling:scheduler}{Round robin} with a \omterm{def:g12:processes-and-scheduling:scheduler}{quantum} of $2$ on the workload of the listing. A, B, C and D arrive at times $0$ to $3$.
Each \omterm{def:g12:processes-and-scheduling:process}{process} runs at most two units, then goes to the end of the queue; B needs a third unit at time $10$, A its
fifth at time $13$; C, the longest, finishes last at $18$.}
\end{omfigure}

The choice of the \omterm{def:g12:processes-and-scheduling:scheduler}{quantum} is a compromise. A short \omterm{def:g12:processes-and-scheduling:scheduler}{quantum} gives each \omterm{def:g12:processes-and-scheduling:process}{process} a turn soon, which is what a person at
a keyboard notices; but each switch costs time. A long \omterm{def:g12:processes-and-scheduling:scheduler}{quantum} wastes little in switches, but \omterm{def:g12:processes-and-scheduling:scheduler}{round robin} then
behaves like first come, first served. The problem of the chapter measures this compromise.

\section{Two processes and one variable}\label{sec:g12:processes-and-scheduling:race}

Processes that share data, such as two threads of a program updating the same counter, can be interrupted between
any two instructions. The line \code{counter += 1} is three instructions for the processor: \emph{load} the counter
into a register, \emph{add} $1$ to the register, \emph{store} the register back. The program below plays two such
processes in every possible order.

\omcode{code/grade-12/10-processes-and-scheduling/listings/interleave.py}{1}{35}{Every interleaving of two
three-step updates.}\label{lst:g12:processes-and-scheduling:interleave}

\omoutput{code/grade-12/10-processes-and-scheduling/outputs/interleave.txt}{Only two orders of the twenty give the
right answer.}\label{out:g12:processes-and-scheduling:interleave}

\begin{omfigure}
\begin{tikzpicture}[font=\footnotesize, x=1.9cm, y=-0.9cm]
  \node[left] at (0.4,0) {P}; \node[left] at (0.4,1) {counter}; \node[left] at (0.4,2) {Q};
  \draw[black!30] (0.5,0) -- (6.6,0); \draw[black!30] (0.5,2) -- (6.6,2);
  \foreach \k/\w/\a/\v/\f in {1/0/load/0/omDef!12, 2/0/add/0/omDef!12, 3/2/load/0/omMeth!15,
      4/0/store/1/omDef!12, 5/2/add/1/omMeth!15, 6/2/store/1/omMeth!15} {
    \node[draw, rounded corners=1pt, fill=\f] at (\k,\w) {\a};
    \node[draw=black!40, fill=black!5] at (\k,1) {$\v$};
  }
  \node[font=\scriptsize, below=1pt] at (1,0.25) {reg $0$};
  \node[font=\scriptsize, below=1pt] at (2,0.25) {reg $1$};
  \node[font=\scriptsize, above=1pt] at (3,1.75) {reg $0$};
  \node[font=\scriptsize, above=1pt] at (5,1.75) {reg $1$};
  \draw[->, thick, omThm] (4,0.3) -- (4,0.7);
  \draw[->, thick, omThm] (6,1.7) -- (6,1.3);
  \node[font=\scriptsize, omThm, right] at (6.15,1) {lost: $1$, not $2$};
\end{tikzpicture}

{\small The interleaving printed above. Q loads the counter before P has stored its result, so both start from $0$;
P stores $1$, then Q stores $1$ over it. One of the two additions is lost.}
\end{omfigure}

\begin{definition}[Race condition, critical section, mutual exclusion, lock]\label{def:g12:processes-and-scheduling:race}
A \emph{race condition}\index{race condition} is a bug whose effect depends on the order in which processes happen to
run. A \emph{critical section}\index{critical section} is a part of a program that uses shared data and must not run
in two processes at once. \emph{Mutual exclusion}\index{mutual exclusion} is the guarantee that at most one \omterm{def:g12:processes-and-scheduling:process}{process} is
in a critical section at any time. A \emph{lock}\index{lock} provides it: a \omterm{def:g12:processes-and-scheduling:process}{process} \emph{acquires} the lock before the
critical section and \emph{releases} it after; a \omterm{def:g12:processes-and-scheduling:process}{process} that tries to acquire a lock already held waits until it is
released.
\end{definition}

\omcode{code/grade-12/10-processes-and-scheduling/listings/with_lock.py}{1}{40}{The two updates inside a \omterm{def:g12:processes-and-scheduling:race}{lock}, with
every possible run.}\label{lst:g12:processes-and-scheduling:lock}

\omoutput{code/grade-12/10-processes-and-scheduling/outputs/with_lock.txt}{Whoever acquires the \omterm{def:g12:processes-and-scheduling:race}{lock} first does its
three steps alone: the counter is always $2$.}\label{out:g12:processes-and-scheduling:lock}

A real \omterm{def:g12:processes-and-scheduling:race}{lock} must itself be acquired without a race: two processes testing ``is it free?'' and then taking it would
repeat the bug. Processors provide special instructions that test and set a memory cell in one indivisible step, and
operating systems build \omterm{def:g12:processes-and-scheduling:race}{locks} on them.

\section{Deadlock}\label{sec:g12:processes-and-scheduling:deadlock}

\begin{definition}[Deadlock]\label{def:g12:processes-and-scheduling:deadlock}
A \emph{deadlock}\index{deadlock} is a state in which each \omterm{def:g12:processes-and-scheduling:process}{process} of a group waits for a resource held by another
\omterm{def:g12:processes-and-scheduling:process}{process} of the group, so that none of them can ever go on.
\end{definition}

The simplest case needs two \omterm{def:g12:processes-and-scheduling:race}{locks}. P takes A, then B; Q takes B, then A. If each takes its first \omterm{def:g12:processes-and-scheduling:race}{lock} before the other
takes its second, both wait for ever.

\omcode{code/grade-12/10-processes-and-scheduling/listings/deadlock_sim.py}{1}{29}{Every run of two processes taking
two \omterm{def:g12:processes-and-scheduling:race}{locks}.}\label{lst:g12:processes-and-scheduling:deadlock}

\omcode{code/grade-12/10-processes-and-scheduling/listings/deadlock_sim.py}{31}{51}{The cycle of the wait-for \omterm{def:g10:graphs-first-look:graph}{graph},
and the two orders of taking the \omterm{def:g12:processes-and-scheduling:race}{locks}.}\label{lst:g12:processes-and-scheduling:deadlock-run}

\omoutput{code/grade-12/10-processes-and-scheduling/outputs/deadlock_sim.txt}{Opposite orders can \omterm{def:g12:processes-and-scheduling:deadlock}{deadlock}; the same
order never does.}\label{out:g12:processes-and-scheduling:deadlock}

\begin{omfigure}
\begin{tikzpicture}[font=\footnotesize, pr/.style={draw, circle, minimum size=10mm, fill=omDef!12},
    rs/.style={draw, minimum size=9mm, fill=omMeth!15}, ed/.style={->, thick, >={Stealth}}]
  \node[pr] (P) at (0,0) {P}; \node[pr] (Q) at (4,0) {Q};
  \node[rs] (A) at (2,1.6) {A}; \node[rs] (B) at (2,-1.6) {B};
  \draw[ed] (A) -- node[above left, font=\scriptsize] {held by} (P);
  \draw[ed, omThm] (P) -- node[below left, font=\scriptsize] {waits for} (B);
  \draw[ed] (B) -- node[below right, font=\scriptsize] {held by} (Q);
  \draw[ed, omThm] (Q) -- node[above right, font=\scriptsize] {waits for} (A);
\end{tikzpicture}

{\small The resource-allocation \omterm{def:g10:graphs-first-look:graph}{graph} of the \omterm{def:g12:processes-and-scheduling:deadlock}{deadlock}. Processes are circles, \omterm{def:g12:processes-and-scheduling:race}{locks} squares; an arrow from a \omterm{def:g12:processes-and-scheduling:race}{lock} to
a \omterm{def:g12:processes-and-scheduling:process}{process} means ``held by'', from a \omterm{def:g12:processes-and-scheduling:process}{process} to a \omterm{def:g12:processes-and-scheduling:race}{lock} ``waits for''. The four arrows form a cycle. Leaving out the
\omterm{def:g12:processes-and-scheduling:race}{locks} gives the wait-for \omterm{def:g10:graphs-first-look:graph}{graph} of the output: P waits for Q, Q waits for P.}
\end{omfigure}

A \omterm{def:g12:processes-and-scheduling:deadlock}{deadlock} needs four conditions together: \emph{\omterm{def:g12:processes-and-scheduling:race}{mutual exclusion}} (a resource is held by one \omterm{def:g12:processes-and-scheduling:process}{process} at a time),
\emph{hold and wait} (a \omterm{def:g12:processes-and-scheduling:process}{process} keeps what it holds while it waits for more), \emph{no preemption} (a resource cannot
be taken away from the \omterm{def:g12:processes-and-scheduling:process}{process} holding it), and \emph{circular wait} (a cycle of processes, each waiting for the
next). Breaking any one prevents \omterm{def:g12:processes-and-scheduling:deadlock}{deadlock}. The usual remedy breaks the fourth: number the resources, and have every
\omterm{def:g12:processes-and-scheduling:process}{process} take them in increasing order. A cycle would then need a \omterm{def:g12:processes-and-scheduling:process}{process} waiting for a resource numbered lower than
one it holds, which the rule forbids; the last line of the output checks it on every run.

\begin{history}[Coffman, Elphick and Shoshani, 1971]
The four conditions were stated together by Edward G.\ Coffman, Michael Elphick and Arie Shoshani in ``System
\omterm{def:g12:processes-and-scheduling:deadlock}{deadlocks}'', a survey published in 1971 in \emph{ACM Computing Surveys}.
% ledger: g12:coffman-1971
\end{history}

\section{Exercises}\label{sec:g12:processes-and-scheduling:exercises}

\begin{exercise}[$\star$]\label{exo:g12:processes-and-scheduling:1}
Three processes arrive at time $0$ in the order X (burst $4$), Y ($1$), Z ($3$). Give the order and the slices of first
come, first served.
\end{exercise}

\begin{exercise}[$\star$]\label{exo:g12:processes-and-scheduling:2}
Compute the waiting time of each \omterm{def:g12:processes-and-scheduling:process}{process} of exercise \ref{exo:g12:processes-and-scheduling:1}, and the mean.
\end{exercise}

\begin{exercise}[$\star$]\label{exo:g12:processes-and-scheduling:3}
Write the Gantt line of \omterm{def:g12:processes-and-scheduling:scheduler}{round robin} with a \omterm{def:g12:processes-and-scheduling:scheduler}{quantum} of $1$ on X, Y, Z of exercise
\ref{exo:g12:processes-and-scheduling:1}.
\end{exercise}

\begin{exercise}[$\star$]\label{exo:g12:processes-and-scheduling:4}
How many \omterm{def:g12:processes-and-scheduling:context-switch}{context switches} does \omterm{def:g12:processes-and-scheduling:scheduler}{round robin} with $q = 2$ make on the workload of the chapter? And with $q = 4$?
\end{exercise}

\begin{exercise}[$\star$]\label{exo:g12:processes-and-scheduling:5}
Write \code{turnaround(workload, slices)}, which returns the turnaround of each \omterm{def:g12:processes-and-scheduling:process}{process} as a dictionary.
\end{exercise}

\begin{exercise}[$\star\star$]\label{exo:g12:processes-and-scheduling:6}
A \omterm{def:g12:processes-and-scheduling:scheduler}{round robin} forgets to put an unfinished \omterm{def:g12:processes-and-scheduling:process}{process} back in the queue. What happens on the workload of the chapter
with $q = 2$?
\end{exercise}

\begin{exercise}[$\star\star$]\label{exo:g12:processes-and-scheduling:7}
How many interleavings do two processes of $3$ steps each have? Of $a$ and $b$ steps? Explain the count.
\end{exercise}

\begin{exercise}[$\star\star$]\label{exo:g12:processes-and-scheduling:8}
Write \code{sjf(workload)}, shortest job first: when the processor is free, it runs to its end the ready \omterm{def:g12:processes-and-scheduling:process}{process} with
the smallest burst. Apply it to the workload of the chapter.
\end{exercise}

\begin{exercise}[$\star\star$]\label{exo:g12:processes-and-scheduling:9}
A function acquires a \omterm{def:g12:processes-and-scheduling:race}{lock}, and returns early, without releasing it, when the value it reads is negative. What
happens to the other processes that use the \omterm{def:g12:processes-and-scheduling:race}{lock}?
\end{exercise}

\begin{exercise}[$\star\star$]\label{exo:g12:processes-and-scheduling:10}
Give the final counter for the order: P load, Q load, Q add, Q store, P add, P store. Is an update lost?
\end{exercise}

\begin{exercise}[$\star\star$]\label{exo:g12:processes-and-scheduling:11}
Two designs. In the first, processes never use the printer: they hand whole documents to a printer server, which
prints them one after the other. In the second, a \omterm{def:g12:processes-and-scheduling:process}{process} asks for all the resources it will need in one request,
granted all at once or not at all, before it starts. Which of the four conditions of \omterm{def:g12:processes-and-scheduling:deadlock}{deadlock} does each design rule
out?
\end{exercise}

\begin{exercise}[$\star\star\star$]\label{exo:g12:processes-and-scheduling:12}
Prove that with the \omterm{def:g12:processes-and-scheduling:race}{lock}, no run of \code{with_lock.py} loses an update, whatever the order of the steps.
\end{exercise}

\begin{exercise}[$\star\star\star$]\label{exo:g12:processes-and-scheduling:13}
Write \code{has_cycle(waits_for)}, which detects a \omterm{def:g12:processes-and-scheduling:deadlock}{deadlock} in a wait-for \omterm{def:g10:graphs-first-look:graph}{graph} given as a dictionary from each
\omterm{def:g12:processes-and-scheduling:process}{process} to the list of processes it waits for.
\end{exercise}

\begin{exercise}[$\star\star\star$]\label{exo:g12:processes-and-scheduling:14}
Four processes arrive at time $0$ with bursts $6$, $8$, $3$, $2$. Compute the mean waiting time of first come, first
served and of shortest job first. Show that, when all processes arrive together, shortest job first gives the smallest
mean waiting time of all orders.
\end{exercise}

\begin{exercise}[$\star\star\star$]\label{exo:g12:processes-and-scheduling:15}
Five philosophers sit at a round table, with a fork between each two. To eat, a philosopher takes the fork on the left,
then the one on the right. Show that a \omterm{def:g12:processes-and-scheduling:deadlock}{deadlock} can happen, and that taking the lower-numbered fork first prevents it.
\end{exercise}

\section{Problem: Choosing the Quantum}\label{sec:g12:processes-and-scheduling:problem}

\begin{problem}[{Weekend problem --- how long should a turn last?}]\label{pb:g12:processes-and-scheduling:1}
Five processes share one core: A, a long computation, arrives first; four short ones follow. Every \omterm{def:g12:processes-and-scheduling:context-switch}{context switch} costs
one unit of time. The program tries \omterm{def:g12:processes-and-scheduling:scheduler}{round robin} with every \omterm{def:g12:processes-and-scheduling:scheduler}{quantum} from $1$ to $12$.

\omcode{code/grade-12/10-processes-and-scheduling/listings/quantum.py}{1}{28}{\omterm{def:g12:processes-and-scheduling:scheduler}{Round robin} on one workload, for every
\omterm{def:g12:processes-and-scheduling:scheduler}{quantum}.}\label{lst:g12:processes-and-scheduling:quantum}

\omoutput{code/grade-12/10-processes-and-scheduling/outputs/quantum.txt}{Mean turnaround, number of switches, and the
time when everything is done.}\label{out:g12:processes-and-scheduling:quantum}

\needspace{4\baselineskip}\medskip\noindent\textbf{Part I --- The workload.}
\begin{enumerate}
  \item How much processor time does the workload need in all?
  \item Why can no schedule finish everything before time $35$?
  \item With first come, first served, in which order do the processes run? Give the slices, with the switches.
  \item Compute the turnaround of each \omterm{def:g12:processes-and-scheduling:process}{process} under first come, first served, and check the mean printed.
\end{enumerate}

\needspace{4\baselineskip}\medskip\noindent\textbf{Part II --- Short and long quanta.}
\begin{enumerate}[resume]
  \item Why is \omterm{def:g12:processes-and-scheduling:scheduler}{round robin} with $q = 1$ worse than first come, first served here?
  \item Count the time lost in switches for $q = 1$, and check it against the time when everything is done.
  \item Why is the result the same for every $q \ge 20$? Which schedule is it?
  \item Why do $q = 12$ and first come, first served differ, although A runs first in both?
\end{enumerate}

\needspace{4\baselineskip}\medskip\noindent\textbf{Part III --- The best quantum.}
\begin{enumerate}[resume]
  \item Which \omterm{def:g12:processes-and-scheduling:scheduler}{quantum} gives the smallest mean turnaround? By how much does it beat first come, first served?
  \item With $q = 6$, write the slices, and find when each \omterm{def:g12:processes-and-scheduling:process}{process} finishes.
  \item Explain why $q = 6$ suits this workload: compare it with the bursts of the short processes.
  \item Why does the mean grow again for $q$ from $7$ to $12$?
  \item The mean does not change steadily with $q$: $q = 5$ is worse than $q = 4$. Compare the finishing time of each
    \omterm{def:g12:processes-and-scheduling:process}{process} for the two quanta, and explain.
\end{enumerate}

\needspace{4\baselineskip}\medskip\noindent\textbf{Part IV --- Other measures.}
\begin{enumerate}[resume]
  \item Compute the response time of B under first come, first served and under $q = 6$.
  \item A person typing in an editor cares about the response time; a person waiting for a report, about the
    turnaround. Which \omterm{def:g12:processes-and-scheduling:scheduler}{quantum} would each prefer here?
  \item If a switch cost nothing, the program gives a mean turnaround of $17.4$ for $q = 1$, and $17.0$ for $q = 2$
    and $q = 3$. Why does $q = 1$ improve so much, and why is it still not the best?
  \item Shortest job first is not \omterm{def:g12:processes-and-scheduling:scheduler}{round robin}. Would it beat the best \omterm{def:g12:processes-and-scheduling:scheduler}{quantum} on this workload? Why can a real
    \omterm{def:g12:processes-and-scheduling:scheduler}{scheduler} not use it directly?
  \item Suppose a \omterm{def:g12:processes-and-scheduling:scheduler}{quantum} of $4$ milliseconds and a switch costing $4$ microseconds. What fraction of the processor
    time goes to switches when every \omterm{def:g12:processes-and-scheduling:scheduler}{quantum} is used in full? Compare with $q = 6$ here.
  \item Why do real systems also change the \omterm{def:g12:processes-and-scheduling:scheduler}{quantum}, or the order, for processes that often block?
  \item State the result: the best \omterm{def:g12:processes-and-scheduling:scheduler}{quantum} for this workload and switch cost, its mean turnaround, and that of first
    come, first served.
\end{enumerate}
\end{problem}
